\documentclass[tikz,border=2pt]{standalone}
\usepackage{ctex}

\begin{document}
% 按原图“图 1.19 考虑自旋–轨道耦合后的核子能级”的能级顺序重绘。
\begin{tikzpicture}[
    >=latex, x=1cm, y=.44cm, font=\small,
    splitting/.style={densely dashed, black!58, line width=.4pt},
    level/.style={line width=.7pt},
    count/.style={minimum width=2.35em, inner sep=1pt},
    closed/.style={count, draw, rounded corners=1pt, font=\small\bfseries}
]
    \node at (1.15,29) {振子激发能};
    \node at (3.1,29) {耦合前 $n\ell$};
    \node at (7.05,29) {耦合后 $n\ell_j$};
    \node at (8.95,29) {容量 $2j + 1$};
    \node at (10.75,29) {累计数};
    \draw[->] (-.05,0) -- (-.05,27.9) node[above] {$E$};

    % 左列按振子壳层分组。短括号涵盖同一壳层内所画的轨道。
    \foreach \shell/\height in {0/0,1/2.35,2/5.6,3/10.65,4/16.3,5/22.45,6/27.55}
        \node at (1.15,\height) {$\shell\hbar\omega$};
    \foreach \low/\high in {5.15/6.05,9.8/11.5,14.35/18.25,20.15/24.75}
        \draw (2.37,\low) -- (2.23,\low) -- (2.23,\high) -- (2.37,\high);

    \foreach \orbital/\height in {
        1s/0, 1p/2.35, 1d/5.15, 2s/6.05, 1f/9.8, 2p/11.5,
        1g/14.35, 2d/16.85, 3s/18.25, 1h/20.15, 2f/22.65,
        3p/24.75, 1i/27.55}
    {
        \node at (2.95,\height) {$\mathrm{\orbital}$};
        \draw[level] (3.35,\height) -- (4.05,\height) coordinate (parent-\orbital);
    }

    % 依能量递增列出：轨道、2j、图上高度、累计粒子数、闭壳层标记。
    \foreach \orbital/\twoj/\height/\total/\style in {
        1s/1/0/2/closed,
        1p/3/1.85/6/count,    1p/1/2.85/8/closed,
        1d/5/4.7/14/count,    2s/1/5.7/16/count,    1d/3/6.7/20/closed,
        1f/7/8.55/28/closed,
        2p/3/10.4/32/count,   1f/5/11.4/38/count,   2p/1/12.4/40/count,
        1g/9/13.4/50/closed,
        1g/7/15.25/58/count,  2d/5/16.25/64/count,  2d/3/17.25/68/count,
        3s/1/18.25/70/count,  1h/11/19.25/82/closed,
        1h/9/21.1/92/count,   2f/7/22.1/100/count,  2f/5/23.1/106/count,
        3p/3/24.1/110/count,  3p/1/25.1/112/count,  1i/13/26.1/126/closed}
    {
        \draw[splitting] (parent-\orbital) -- (5.5,\height);
        \draw[level] (5.5,\height) -- (6.48,\height);
        \node[anchor=west] at (6.52,\height) {$\mathrm{\orbital}_{\twoj/2}$};
        \node at (8.95,\height) {\the\numexpr\twoj+1\relax};
        \node[\style] at (10.75,\height) {\total};
    }
    \label{spin-orbit-splitting}
\end{tikzpicture}
\end{document}
